\documentclass[12pt, a4paper]{academics}

\academicsCoverLabel{解题报告}
\academicsCourse{算法设计与分析}{Design and Analysis of Algorithms}
\academicsReport{字符串匹配}{2026 年 6 月 8 日}
\academicsArthur{华博文}{19240212}

\begin{document}

\makeacademicscover

\section{引言}

字符串匹配是计算机科学中最基础也最具实用价值的问题之一——给定文本串与模式串，找出模式串在文本中的所有出现位置。从文本编辑器的查找替换到生物信息学中的基因序列比对，从网络入侵检测的签名匹配到搜索引擎的倒排索引构建，字符串匹配算法渗透在计算机系统的各个层面，其效率往往直接决定了上层应用的响应速度。

本报告以字符串匹配为主线，依次展开四道关联紧密的题目。从最朴素的滑动窗口匹配开始，引入 Knuth-Morris-Pratt（以下简称 KMP）算法，利用前缀函数的失配跳转机制将复杂度从 $\mathrm{O}(nm)$ 压缩至 $\mathrm{O}(n + m)$；继而将视角转向回文子串问题，探讨 Manacher 算法如何通过对称性复用实现线性时间的最长回文子串求解；最后进入多模式匹配领域，展示 Aho-Corasick（以下简称 AC）自动机如何将 KMP 的单模式失配跳转思想推广到基于 trie 的多模式同时匹配场景。四道题目的核心思想一脉相承——利用已经获得的匹配信息，避免重复比较。

报告中的所有代码均使用 C++23 标准编写，在 Apple clang 21.0.0 编译器下通过测试。

\section{题一：朴素字符串匹配}

\subsection{问题重述}

给定一个文本串 $S$（长度为 $n$）和一个模式串 $P$（长度为 $m$），字符串仅由小写英文字母组成，要求找出 $P$ 在 $S$ 中所有出现的起始位置（从 $0$ 开始编号），按升序输出。若 $P$ 未在 $S$ 中出现，输出 $-1$。本题的数据范围为 $1 \le n, m \le 1000$，允许 $\mathrm{O}(nm)$ 的朴素算法。

\subsection{滑动窗口与穷举验证}

解决这一问题的最直接思路是穷举所有可能的匹配起始位置并逐一验证。等价地，我们可以想象一个长度为 $m$ 的窗口在文本串 $S$ 上从最左端滑向最右端，在每个窗口位置检查窗口内的子串是否与模式串 $P$ 完全一致。令窗口左端点为 $i$（$0 \le i \le n - m$），逐字符比较 $S[i \dots i+m-1]$ 与 $P[0 \dots m-1]$，一旦发现某个字符不匹配，立即跳出内层比较并尝试下一个起始位置。

这一算法的正确性是不言自明的——它穷举了所有 $n - m + 1$ 个候选起始位置，并对每一个候选位置进行了完整的验证，不会遗漏任何匹配。算法的主体是两重嵌套循环：外层遍历起始位置，内层执行逐字符比对。

在复杂度方面，外层循环遍历 $\mathrm{O}(n)$ 个候选起始位置。在最坏情况下（例如 $S$ 和 $P$ 均由同一个字符组成，或者均形如 $S = \text{aaa}\dots\text{ab}$、$P = \text{aaa}\dots\text{a}$），内层循环需要匹配全部 $m$ 个字符才能确认一次匹配，总时间复杂度为 $\mathrm{O}(nm)$。最好情况下，每个起始位置的第一个字符即失配，复杂度为 $\mathrm{O}(n)$。空间上仅需存储输入字符串和结果数组，空间复杂度为 $\mathrm{O}(n + m)$。对于 $n, m \le 1000$ 的数据规模，至多约 $10^6$ 次字符比较在毫秒级别内即可完成。

然而，当数据范围扩展至 $n, m \le 10^6$ 时，$\mathrm{O}(nm)$ 将产生 $10^{12}$ 级别的运算量，远远超出时限。这意味着我们必须摆脱穷举的思维框架，寻找更高效的匹配策略——这正是下一道题目的出发点。

\subsection{代码实现}

\inputminted{c++}{../src/p1_naive.cpp}

\section{题二：KMP 算法}

\subsection{问题重述}

本题的问题描述与题一完全一致，唯一的变化在于数据范围大幅提升至 $1 \le n, m \le 10^6$。$\mathrm{O}(nm)$ 的朴素算法在此规模下已完全不可行，必须设计线性时间算法。

\subsection{前缀函数与失配跳转}

朴素算法效率低下的根本原因在于：每次失配后，模式串回退到起始位置，文本串也回退到下一个起始位置，此前已经成功匹配的那一部分字符信息被彻底丢弃。例如，当模式串在位置 $j$ 与文本串失配时，我们已知 $S[i-j \dots i-1] = P[0 \dots j-1]$，其中蕴含了模式串自身的结构信息——如果 $P[0 \dots j-1]$ 的某个真前缀恰好等于它的某个真后缀，那么当我们试图从下一个可能的位置重新匹配时，这个公共部分无需再次比较。朴素算法却无视了这一信息，反复扫描已经确认相等的字符，造成了大量冗余。

Knuth、Morris 和 Pratt 在 1977 年各自独立提出的算法正是针对这一冗余的解决方案。其核心数据结构是 next 数组（或称前缀函数 $\pi$）：对于模式串的每个前缀 $P[0 \dots i]$，$\mathrm{next}[i]$ 定义为该前缀的最长真前缀的长度，满足该真前缀同时也是该前缀的后缀。形式化地，

\[
\mathrm{next}[i] = \max\{ k \mid k < i + 1 \text{ 且 } P[0 \dots k-1] = P[i-k+1 \dots i] \}
\]

拥有了 next 数组之后，匹配过程不再需要回溯文本指针。当 $S[i] \neq P[j]$ 时，我们将 $j$ 跳转到 $\mathrm{next}[j-1]$——即已匹配部分的最长公共前后缀长度——然后保持 $i$ 不动，继续比较 $S[i]$ 与新的 $P[j]$。这背后的逻辑是：既然 $S[i-j \dots i-1] = P[0 \dots j-1]$，而已匹配前缀 $P[0 \dots j-1]$ 的最长公共前后缀长度为 $\mathrm{next}[j-1]$，那么模式串的前 $\mathrm{next}[j-1]$ 个字符已经天然地与文本串的末尾 $\mathrm{next}[j-1]$ 个字符对齐，可以直接从模式串的第 $\mathrm{next}[j-1]$ 个位置继续比较。文本指针 $i$ 始终单向推进，每个文本字符至多被比较常数次。

\subsection{两阶段递推与线性复杂度}

构建 next 数组本身是一个递推的自我匹配过程：将模式串对自身运行 KMP 匹配逻辑。从 $i = 1$ 开始，维护 $j$ 表示当前已匹配的真前缀长度。当 $P[i] \neq P[j]$ 时，反复将 $j$ 回退到 $\mathrm{next}[j-1]$ 直到匹配或 $j = 0$；若 $P[i] = P[j]$，则 $j$ 加一，然后将 $\mathrm{next}[i]$ 设为 $j$。

这一构建过程的复杂度为 $\mathrm{O}(m)$，原因在于指针 $j$ 每次至多增加 $1$（总共增加不超过 $m$ 次），而每次回退至少减少 $j$ 的值，因此总的回退次数也被 $m$ 所限制。匹配阶段同理：文本指针 $i$ 从 $0$ 到 $n-1$ 单向移动，模式指针 $j$ 的增减与构建阶段一致，总复杂度为 $\mathrm{O}(n)$。因此，KMP 算法的总体时间复杂度为 $\mathrm{O}(n + m)$，空间复杂度为 $\mathrm{O}(m)$（仅需存储 next 数组）。当 $n$ 和 $m$ 达到 $10^6$ 级别时，线性复杂度使得算法在毫秒级别内给出结果。

从方法论的角度看，KMP 算法的精妙之处在于它将失配时「下一个可能匹配的位置」编码为一个预先计算的数组，从而将冗余回溯替换为常数时间的数组查表。这一「预处理 + 在线查询」的两阶段模式，是字符串算法乃至整个算法设计领域中反复出现的范式。

\subsection{代码实现}

\inputminted{c++}{../src/p2_kmp.cpp}

\section{题三：Manacher 算法求最长回文子串}

\subsection{问题重述}

给定一个长度为 $n$ 的字符串 $S$（$n \le 10^6$，仅含小写英文字母），求其最长回文子串的长度。回文子串是指正读与反读完全相同的连续子串。与题一、题二的不同之处在于，本题并非在两个字符串之间寻找匹配，而是在单一字符串内部寻找具有对称性质的子结构。

\subsection{统一奇偶与对称性复用}

回文子串问题表面上看似与字符串匹配无关，但两者在算法思想上有着深刻的共鸣——都涉及利用已获得的结构信息来跳过重复检查。在 KMP 中，我们利用已匹配前缀的真前后缀相等性来跳过无望的起始位置；在 Manacher 算法中，我们利用已发现回文的对称性来直接复用其镜像位置的半径信息。

Manacher 算法的第一步是一个巧妙的预处理：在原字符串每两个字符之间（以及首尾）插入一个特殊分隔符，再在两端添加哨兵字符。例如字符串 ababa 被变换为 \mil{\^{}\#a\#b\#a\#b\#a\#\$}。这一变换的精妙之处在于，它将奇长度和偶长度的回文统一为同一类问题——变换后的字符串中，每个回文都以某个字符为中心且长度为奇数，原始字符串中的回文长度恰好等于变换后对应回文的半径。这种「统一奇偶」的技巧避免了对奇回文和偶回文分别编写中心扩展逻辑的繁琐。

算法的核心在于维护两个变量：当前已知的最靠右回文的中心 $\mathit{center}$ 和其右边界 $\mathit{right}$。处理到位置 $i$ 时，其关于 $\mathit{center}$ 的对称点为 $\mathit{mirror} = 2 \times \mathit{center} - i$。若 $i < \mathit{right}$，即 $i$ 落在已知最靠右回文的内部，则对称性保证了以 $i$ 为中心的回文在以 $\mathit{center}$ 为中心的回文内部的部分，必然与以 $\mathit{mirror}$ 为中心的回文对称相等。因此 $i$ 的初始回文半径可以直接取 $\min(\mathit{right} - i, p[\mathit{mirror}])$——前者保证不超出已知边界，后者则是镜像处的已知半径。在这两种限制中取较小者，得到了无需逐字符比较即可确定的半径下界。

\subsection{中心扩展与边界维护}

获得初始半径后，算法尝试从这个下界出发向外逐字符扩展，直到遇到不匹配的字符。若扩展后 $i$ 的右边界 $i + p[i]$ 超越了当前的 $\mathit{right}$，则更新 $\mathit{center}$ 和 $\mathit{right}$ 为 $i$ 和 $i + p[i]$——这意味着我们发现了更靠右的回文，其内部的对称信息将在后续位置的处理中被复用。

这一设计的效率保证来自对 $\mathit{right}$ 边界移动次数的分析：$\mathit{right}$ 单调不减，在遍历过程中至多从 $0$ 增加到 $2n$。每次成功扩展都会将 $\mathit{right}$ 向右推进至少一，因此所有位置上的成功扩展总次数不超过 $2n$。每次失败的扩展尝试（边界外第一个字符不匹配）在每个位置上至多发生一次，因此总计 $\mathrm{O}(n)$。由此，算法的总体时间复杂度为 $\mathrm{O}(n)$。空间上需要额外存储变换后的字符串和半径数组，空间复杂度同为 $\mathrm{O}(n)$。

对称性复用这一技巧——利用已探索区域的结构信息为未探索区域提供初始下界，再在此基础上进行有限扩展——不仅限于回文问题。在计算几何的最近点对、序列比对的 banded DP 等场景中，类似的「以已知约束未知」策略均发挥着关键作用。

\subsection{代码实现}

\inputminted{c++}{../src/p3_manacher.cpp}

\section{题四：Aho-Corasick 自动机}

\subsection{问题重述}

给定 $k$ 个模式串 $P_1, P_2, \dots, P_k$（模式串总长度 $\sum |P_i| \le 10^5$，仅含小写英文字母）和一个文本串 $S$（$|S| \le 10^6$），要求对于每个模式串 $P_i$，输出它在 $S$ 中出现的次数，允许重叠匹配。本题是题二 KMP 算法从单模式匹配到多模式匹配的直接推广。

\subsection{从单模式前缀函数到多模式 fail 指针}

KMP 为单模式串构建一维的 next 数组，将失配时回退的位置编码为前缀函数；AC 自动机则为多个模式串构建一棵 trie（前缀树），并在 trie 的每个节点上引入 fail 指针——指向该节点所代表字符串的最长真后缀，且该后缀恰好是某个模式串的前缀。两者在结构上完全对应：next 数组的自我匹配递推对应于 fail 指针在 trie 上的 BFS 传播，失配时的跳转操作对应于沿 fail 链向上爬升。区别仅在于 KMP 面对的是长度为 $m$ 的一维序列，而 AC 自动机面对的是由多个模式串编织而成的树状结构。

算法分为三个清晰的阶段。第一阶段，将所有模式串插入 trie：从根节点出发，沿每个字符的对应边向下走，若边不存在则创建新节点；到达模式串末尾时，在终止节点记录该模式的索引（同一节点可能对应多个模式串，当一个模式串是另一个模式串的后缀时便会出现此种情况）。第二阶段，通过广度优先搜索（BFS）逐层建立 fail 指针。根节点的直接子节点的 fail 指针均指向根节点，因为它们的单字符前缀没有任何真后缀能同时是前缀。对于一般情况：设节点 $u$ 沿字符 $c$ 到达子节点 $v$，则 $v$ 的 fail 指针应从 $u$ 的 fail 指针出发，沿字符 $c$ 能到达的节点；若 $u$ 的 fail 没有对应的字符转移，则递归地沿 fail 链向上寻找（这一过程在实现中通过循环完成）。第三阶段，将文本串 $S$ 逐字符输入自动机：从根节点开始，沿 trie 边转移，若当前节点没有对应的字符转移则沿 fail 指针回退，直到找到转移或回到根节点。每次完成状态转移后，沿 fail 链向上爬升，收集沿途所有终止节点对应的模式串匹配——因为若当前位置匹配了模式串 abc，而有一个模式串 bc 恰好以同一位置结尾，则它必然位于 abc 的 fail 链上。

\subsection{BFS 逐层传播与文本匹配}

构建 trie 的时间复杂度为 $\mathrm{O}(\sum |P_i|)$，因为每个模式串的每个字符恰好被处理一次。构建 fail 指针的 BFS 过程中，每个节点的 fail 链回退总量被 trie 的节点总数所约束（类似于 KMP 中 $j$ 指针的均摊分析），因此该阶段的时间复杂度同样为 $\mathrm{O}(\sum |P_i|)$。文本匹配阶段，文本指针在 trie 上的转移是单向的 $\mathrm{O}(|S|)$；沿 fail 链回溯收集匹配的部分，在最坏情况下可达 $\mathrm{O}(\sum |P_i| \times |S|)$（例如模式串 a, aa, aaa, $\dots$ 全部匹配且反复遍历长 fail 链），但在本题的数据约束下可以接受。空间复杂度为 $\mathrm{O}(\sum |P_i| \times |\Sigma|)$，其中字母表大小 $|\Sigma| = 26$。

AC 自动机将 KMP 的「一维失配跳转」推广为「trie 上的多维失配跳转」，其方法论意义在于展示了如何将单串上的前缀函数思想系统地扩展到多串共享的前缀结构上。从单串到多串，增加的并非仅仅是问题规模，而是问题结构的维度——单串的前缀关系天然是全序的，而多串的前缀关系构成了一棵偏序的树。fail 指针正是将这种偏序结构重新映射为失配时的线性回退路径，从而保持了与 KMP 相同的均摊效率。

\subsection{代码实现}

\inputminted{c++}{../src/p4_ac.cpp}

\section{总结}

本报告围绕字符串匹配这一经典问题，选取了四道层次递进、思想关联的题目，系统展示了从朴素穷举到线性时间优化再到多模式推广的完整算法设计脉络。

题一的朴素窗口匹配是所有后续讨论的原点。它以 $\mathrm{O}(nm)$ 的双重循环穷举了所有可能的匹配位置，算法直接、正确性自明，是小规模数据下的实用选择，也是理解更高级算法的参照基线。

题二引入的 KMP 算法突破了穷举的框架。其核心突破在于将模式串自身的结构信息——前缀与后缀的相等关系——预先编码为 next 数组，使得失配时能够直接跳转到下一个有意义的位置，避免了对文本指针的回溯。从 $\mathrm{O}(nm)$ 到 $\mathrm{O}(n + m)$ 的跨越，本质上是将重复比较替换为查表操作，体现了「空间换时间」与「预处理加速在线查询」两类最基本的算法优化策略。

题三的 Manacher 算法将信息复用的思想从串间匹配延伸至串内对称性探测。通过统一奇偶的预处理和对称半径的继承机制，算法在遍历过程中不断将已探索回文的对称信息投射到未访问的位置，从而将暴力算法的 $\mathrm{O}(n^3)$ 或朴素中心扩展的 $\mathrm{O}(n^2)$ 压缩至 $\mathrm{O}(n)$。Manacher 与 KMP 共享了同一个方法论内核，但作用的对象从串与串之间的前缀关系切换为串内部的对称关系，显示出这一思想的高度可迁移性。

题四的 AC 自动机完成了从单模式到多模式的维度跨越。它不仅在技术上——通过 trie 上的 fail 指针——忠实地推广了 KMP 的前缀函数，更在方法论上揭示了算法推广的一般路径：将一维链式结构上的递推关系映射到树状结构上的层次传播关系。BFS 逐层构建 fail 指针的过程，与 KMP 中从左到右递推 next 数组的过程在逻辑上完全同构，区别仅在于操作对象的拓扑——一个是线性序列，一个是前缀树。

回顾四道题目的递进脉络，贯穿始终的主线是「记忆与复用」：不将已探索过的信息丢弃，而是将其结构化地编码（KMP 的 next 数组、Manacher 的中心与半径、AC 自动机的 fail 指针），从而在后续步骤中重复利用，避免冗余计算。这一思想超越字符串匹配领域，广泛存在于动态规划的记忆化搜索、并查集的路径压缩、线段树的懒标记、Dijkstra 算法的已确定集合维护等众多经典算法设计中。识别可复用的中间信息、设计高效的数据结构加以存储、在恰当的时机加以利用——这三步构成了算法优化的通用路径，也是算法设计中最值得反复锤炼的核心能力。

\end{document}
